\section{Introduction}
\label{sec:introduction}

Training a neural network across many workers requires every worker to agree on the same gradient after each backward pass. In synchronous data-parallel training, each worker computes a gradient from its own shard of a batch, and all workers then need the elementwise sum (or average) of every worker's gradient before taking an optimizer step. This all-to-all reduction, an allreduce, runs once per training step, and at the scale of modern large-model training it is frequently the dominant cost standing between raw compute and wall-clock training time. How efficiently a system performs allreduce over its interconnect, not how many FLOPs its accelerators can issue, is often the practical ceiling on how fast a large training job can go.

A naive allreduce, gather every worker's contribution to one coordinator, sum it there, and broadcast the result back out, puts all $N$ workers' data through a single node twice. That node's link becomes a bottleneck whose cost grows with $N$, which is exactly backwards: adding more workers should not make the communication step per worker more expensive. Ring-allreduce, the algorithm behind GPU communication libraries such as NVIDIA's NCCL \citep{nvidia2024nccl} and popularized for distributed deep learning by Baidu's and Horovod's ring-allreduce implementations, avoids this. Every worker only ever talks to its two ring neighbors, and the total data any one worker sends or receives is, asymptotically, independent of $N$. That property, bandwidth optimality independent of process count, is what makes it the default building block for large-scale synchronous training.

This project implements ring-allreduce from scratch in C++ directly on top of MPI's non-blocking point-to-point primitives (\texttt{MPI\_Isend}/\texttt{MPI\_Irecv}/\texttt{MPI\_Waitall}), with no MPI collective call anywhere inside the implementation, and benchmarks it against vendor \texttt{MPI\_Allreduce} across a message-size sweep from 8 bytes to 128 MiB and a process-count sweep from 2 to 16 ranks. Section~\ref{sec:background} derives the algorithm's two phases and its step-count relative to alternative collective algorithms. Section~\ref{sec:implementation} describes how the implementation handles the edge cases (single rank, message counts smaller than the process count, non-power-of-two process counts) that are where these implementations typically break. Section~\ref{sec:experimental-setup} documents exactly how the benchmark was run and on what. Section~\ref{sec:results} walks through the measured bus bandwidth, the fitted cost model, and the efficiency and step-count comparisons those measurements support. Section~\ref{sec:discussion} is the part that matters most for judging whether this project's own implementation is any good: it uses the fitted model to separate two distinct small-message bottlenecks, pipeline fill latency and rank synchronization overhead, quantifies each, and explains the specific circumstances under which a correct, bandwidth-optimal from-scratch ring can still lose to Open~MPI's own collective algorithm selection. The key experiment is benchmarking \texttt{MPI\_Allreduce} both as it selects algorithms by default and with it forced, via MCA parameters, to run its own ring: the from-scratch ring is about five times slower than the default selection at 8 bytes, but so is Open~MPI's own ring, which shows that the small-message gap is the price of the ring \emph{algorithm}'s step count rather than a deficiency of this implementation. Section~\ref{sec:conclusion} summarizes what was and was not shown here and what a next iteration of this project would need to add.
